\chapter{Markovian Diffusions and Inventory-Sensitive Quotes}\label{ch:markov}
With the economic objective fixed, the simplest predictive state is the current diffusion price and inventory. We first derive the CARA quoting mechanism, then distinguish additive ABM, multiplicative GBM and correlated multi-asset risk. These formulas provide the benchmark against which memory and state compression will be interpreted.
\section{Solving the Optimal Control Problem}


\subsection{The HJB Equation and Optimal Spreads}\label{sec:classical-hjb}
As established in the previous chapter, the market maker's goal is to find the optimal quoting policies $\delta^{bid}$ and $\delta^{ask}$ that maximize the expected CARA utility of terminal wealth. According to the dynamic programming principle, this stochastic optimal control problem can be characterized by its infinitesimal generator. By applying Itô's Lemma for jump-diffusion processes, the problem translates into the Hamilton-Jacobi-Bellman (HJB) equation.

Assuming for now a generic Markovian dynamics for the mid-price $S_t$ characterized by an instantaneous variance $\sigma^2(t,S_t)=\sigma^2$, the HJB equation for the dealer's value function $u(t,s,q,x)$ takes the form

\begin{align*}
    \frac{\partial u}{\partial t} + \frac{1}{2}\sigma^2 \frac{\partial^2 u}{\partial s^2} &+ \max_{\delta^{bid}} \left\{ \lambda(\delta^{bid}) [u(t, s, q+1, x - s + \delta^{bid}) - u(t,s,q,x)] \right\} \\
    &+ \max_{\delta^{ask}} \left\{ \lambda(\delta^{ask}) [u(t, s, q-1, x + s + \delta^{ask}) - u(t,s,q,x)] \right\} = 0.
\end{align*}

This non-linear partial differential equation expresses a continuous equilibrium condition. The first block captures time decay, the second block captures price-diffusion risk, and the two jump blocks capture the expected utility gain or loss generated by bid and ask executions. In equilibrium, these components must balance under the optimal choice of quote distances. 

A critical feature of the exponential utility function $U(W) = -e^{-\gamma W}$ is that it allows the value function to be wealth-independent. We can factor out the cash component $x$ by applying
\begin{align*}
    u(t,s,q,x) = -e^{-\gamma x}e^{-\gamma \theta(s,q,t)},
\end{align*}
where $\theta(s,q,t)$ represents the certainty equivalent of the dealer's inventory position. Substituting this into the HJB equation and dividing by the non-zero common factor $u$, the cash variable $x$ is eliminated, and the problem reduces to an equation exclusively in terms of $\theta$
\begin{align*}
    \frac{\partial \theta}{\partial t} + \frac{1}{2}\sigma^2 \frac{\partial^2 \theta}{\partial s^2} - \frac{1}{2}\gamma \sigma^2 \left(\frac{\partial \theta}{\partial s}\right)^2 &+ \max_{\delta^{bid}} \frac{\lambda(\delta^{bid})}{\gamma} \left[ 1 - e^{-\gamma (\delta^{bid} - s + \theta(q+1) - \theta(q))} \right] \\
    &+ \max_{\delta^{ask}} \frac{\lambda(\delta^{ask})}{\gamma} \left[ 1 - e^{-\gamma (\delta^{ask} + s + \theta(q-1) - \theta(q))} \right] = 0.
\end{align*}
For compactness, the notation $\theta(q)$, $\theta(q+1)$, and $\theta(q-1)$ in this display suppresses the common dependence on $(s,t)$.

A powerful structural property of this formulation is that the optimization of the quoting distances can be decoupled from the explicit solution of the certainty equivalent $\theta$. To see this, we define the dealer's reservation prices for buying and selling an additional unit of inventory as the marginal change in the certainty equivalent
\begin{align*}
    &r^{bid}(s,q,t) = \theta(s,q+1,t) - \theta(s,q,t),\\
    &r^{ask}(s,q,t) = \theta(s,q,t) - \theta(s,q-1,t).
\end{align*}
Focusing on the bid side, the embedded optimization problem becomes 
\begin{align*}
    \max_{\delta^{bid}} \lambda(\delta^{bid}) [1 - e^{-\gamma(\delta^{bid} - s + r^{bid})}].
\end{align*}
The first-order condition with respect to $\delta^{bid}$ yields an exact, model-free relationship between the optimal spread and the elasticity of the limit order book
\begin{align*}
    s - r^{bid} =\delta^{bid} - \frac{1}{\gamma} \ln \left( 1 - \gamma \frac{\lambda(\delta^{bid})}{\lambda'(\delta^{bid})} \right).
\end{align*}
Under our assumption of exponentially decaying arrival intensities, $\lambda(\delta) = A e^{-k\delta}$, the ratio $\lambda / \lambda'$ simplifies strictly to $-1/k$. Therefore, the optimal quoting distances relative to the indifference prices are derived as
\begin{align*}
    &\delta^{bid} = (s - r^{bid}) + \frac{1}{\gamma} \ln \left( 1 + \frac{\gamma}{k} \right),\\
    &\delta^{ask} = (r^{ask} - s) + \frac{1}{\gamma} \ln \left( 1 + \frac{\gamma}{k} \right).
\end{align*}
This result gives a useful separation of roles. Under exponential fill intensities, the liquidity offset around the reservation prices is governed by the execution-sensitivity parameter $k$ and by risk aversion $\gamma$. The final level of the bid and ask quotes still depends on the reservation prices $r^{bid}$ and $r^{ask}$, and therefore on price risk, inventory, and the remaining horizon.

\subsection{The Certainty Equivalent and Reservation Price}
To make the optimal quoting distances derived in the previous section operational, we must explicitly determine the certainty equivalent $\theta(s,q,t)$ and, consequently, the reservation prices $r^{bid}$ and $r^{ask}$. 

As shown in the reduced HJB equation, the dynamics of $\theta$ depend heavily on the continuous diffusive term, which is driven by the underlying mid-price process. Following the Avellaneda--Stoikov baseline benchmark, we assume the mid-price follows an Arithmetic Brownian Motion (ABM) with constant volatility $\sigma$
\begin{align*}
    dS_t = \sigma dW_t.
\end{align*}
Under this specification, the pure inventory risk penalty can be isolated by solving the HJB equation for a "frozen inventory" portfolio (i.e., ignoring the jump terms associated with spread capture). In this case, the certainty equivalent admits a closed-form solution. The subjective value to the dealer of holding $q$ shares is simply the marked-to-market value of the portfolio minus a variance penalty
\begin{align*}
    \theta(s,q,t) = s \cdot q - \frac{1}{2}\gamma\sigma^2 q^2 (T-t).
\end{align*}
Applying the marginal difference definitions established in Section~\ref{sec:classical-hjb}, we can compute the exact reservation bid and ask prices. The maximum price the dealer is willing to pay to acquire an additional share is
\begin{align*}
    r^{bid}(s,q,t) = \theta(s,q+1,t) - \theta(s,q,t) = s - q\gamma\sigma^2(T-t) - \frac{1}{2}\gamma\sigma^2(T-t).
\end{align*}
Conversely, the minimum price the dealer is willing to accept to sell a share is
\begin{align*}
    r^{ask}(s,q,t) = \theta(s,q,t) - \theta(s,q-1,t) = s - q\gamma\sigma^2(T-t) + \frac{1}{2}\gamma\sigma^2(T-t).
\end{align*}
For practical implementation, it is convenient to define a single, central reservation price $r(s,q,t)$ as the midpoint between the bid and ask reservation levels:
\begin{align*}
    r(s,q,t) = \frac{r^{bid}(s,q,t) + r^{ask}(s,q,t)}{2} = s - q\gamma\sigma^2(T-t).
\end{align*}

This formula represents the dealer's personalized ``fair value'' for the asset and constitutes the core of the inventory control mechanism. It clearly highlights the linear inventory effect under the ABM assumption: if the dealer is long ($q > 0$), the reservation price is shifted below the market mid-price, reflecting an urgency to sell. If short ($q < 0$), the reference price is shifted upwards. 

Crucially, the magnitude of this shift (and thus the strength of the inventory control) is directly proportional to the accumulated inventory $q$, the risk aversion parameter $\gamma$, the constant variance $\sigma^2$, and the time remaining to the end of the trading horizon $(T-t)$.

\subsection{The Inventory-Based Strategy (Avellaneda--Stoikov)}\label{sec:classical-policy}
By combining the results derived in the previous sections, we can finally define the operational market making strategy proposed by Avellaneda and Stoikov. The strategy is characterized by two distinct components: a microstructural component $\delta$, which dictates the width of the spread, and a risk-management component $r$, which dictates its position.

The optimal bid and ask quotes, $p^{bid}$ and $p^{ask}$, are obtained by centering the optimal distance $\delta$ around the central reservation price $r(s,q,t)$
\begin{align*}
    p^{bid}(s,q,t) = r(s,q,t) - \frac{1}{\gamma} \ln \left( 1 + \frac{\gamma}{k} \right),\\
    p^{ask}(s,q,t) = r(s,q,t) + \frac{1}{\gamma} \ln \left( 1 + \frac{\gamma}{k} \right).
\end{align*}
Substituting the closed-form expression for $r$ under the ABM assumption, the strategy can be explicitly written as
\begin{align*}
    &p^{bid} = \underbrace{s - q\gamma\sigma^2(T-t)}_{\text{Reservation Price}} - \underbrace{\frac{1}{\gamma} \ln \left( 1 + \frac{\gamma}{k} \right)}_{\text{Spread Distance } \delta},\\
    &p^{ask} = \underbrace{s - q\gamma\sigma^2(T-t)}_{\text{Reservation Price}} + \underbrace{\frac{1}{\gamma} \ln \left( 1 + \frac{\gamma}{k} \right)}_{\text{Spread Distance } \delta}.
\end{align*}

This formulation reveals the core logic of inventory-based market making: the dealer maintains a constant total spread, but its center is "skewed" relative to the market mid-price $s$ based on the current inventory $q$. 

Specifically, we can identify different behaviors based on the $q$ value. If $q=0$ (neutral inventory), the reservation price equals the mid-price ($r=s$). The dealer quotes symmetrically around the market price, acting as a pure liquidity provider. If $q>0$ (long position), the reservation price shifts downward ($r < s$). Both the bid and ask quotes are lowered and this makes the dealer's ask more competitive (increasing the probability of selling) and the bid less attractive (decreasing the probability of buying more shares), effectively inducing a mean-reverting behavior on the inventory. If $q<0$ (short position), the reservation price shifts upward ($r > s$). The quotes are raised, making the bid more competitive to attract sellers and the ask less attractive to avoid further shorting.

The intensity of this skewing effect is dynamically adjusted according to the remaining time horizon $(T-t)$. In the finite-horizon formula used here, the inventory skew is proportional to $(T-t)$; therefore, absent an additional terminal inventory penalty or a forced liquidation constraint, the skew decreases as $t$ approaches $T$ and the reservation price converges back toward the mid-price. The advantage of the Avellaneda--Stoikov framework over naive quoting strategies is thus the endogenous inventory-sensitive quote displacement generated during the remaining risk horizon.

\subsection{Benchmark Strategies for Simulation}

To systematically evaluate the performance and the risk-mitigation capabilities of the inventory-based policy derived in Section~\ref{sec:classical-policy}, we must establish a clear baseline for comparison. In the simulation experiments conducted in the subsequent chapters, the Avellaneda--Stoikov model will be tested against a primary benchmark: the Naive Symmetric Strategy.

The naive symmetric strategy represents the behavior of a traditional liquidity provider who optimizes for spread capture based on market microstructure, but completely ignores inventory risk. Mathematically, this behavior corresponds to a dealer who calculates the optimal quoting distance $\delta$ using the limit order book elasticity $k$, but blindly anchors these quotes to the current market mid-price $s$ rather than their personal reservation price $r$.

Using the same optimal distance $\delta$ defined previously, the quotes for the symmetric dealer are simply
\begin{align*}
    &p^{bid}_{sym}(s,t) = s - \delta = s - \frac{1}{\gamma} \ln \left( 1 + \frac{\gamma}{k} \right),\\
    &p^{ask}_{sym}(s,t) = s + \delta = s + \frac{1}{\gamma} \ln \left( 1 + \frac{\gamma}{k} \right).
\end{align*}
By keeping the quotes perfectly centered around the mid-price ($p^{bid}_{sym}$ and $p^{ask}_{sym}$ are equidistant from $s$), this strategy maximizes the probability of symmetric executions in a static market. However, because the reference price does not adjust to the accumulated inventory $q$, the symmetric dealer lacks any endogenous mechanism to mean-revert their position. 

If the market exhibits a strong directional trend, the symmetric dealer will continuously execute on one side of the book, accumulating a large directional inventory (e.g., buying in a falling market). Without the "skewing" mechanism to protect the portfolio, the marked-to-market wealth $W_t$ of the symmetric strategy is exposed to severe drawdowns. 

Comparing the inventory-based strategy (which actively shifts $r$) against this symmetric benchmark (which keeps $r=s$) isolates the exact value of the inventory control mechanism. This dynamic will be the central focus of the time-discretized Monte Carlo simulations developed in the following chapters, where both strategies will be subjected to identical order flows and price paths across different diffusion assumptions.

\section{Classical Diffusions under a Finite Horizon}


\subsection{Case 1: Arithmetic Brownian Motion with Finite Horizon}\label{sec:abm}

In this initial benchmark scenario, we model the dynamics of the asset's mid-price through an Arithmetic Brownian Motion (ABM). Assuming a constant drift $\mu$ and a constant volatility parameter $\sigma$, the stochastic differential equation (SDE) governing the price process is defined as
\begin{align*}
    dS_t = \mu dt + \sigma dW_t.
\end{align*}
The decision to employ an ABM, deviating from the Geometric Brownian Motion (GBM) typically used in derivative pricing, is motivated by several structural and economic considerations. Primarily, the ABM provides the most transparent setting to isolate and interpret inventory risk, since price increments are additive, the marked-to-market effect of holding a position is strictly linear with respect to the inventory size. Furthermore, this price dynamics offers superior analytical tractability when paired with the Constant Absolute Risk Aversion (CARA) utility function. Under the ABM assumption, terminal wealth is normally distributed, allowing the moment-generating function to factorize and yielding elegant closed-form solutions for the dealer's reservation price. This choice also ensures that the utility functionals remain bounded throughout the optimization process.

The optimization in this case is conducted over a fixed time horizon $[0, T]$, which represents a standard trading session. Accounting for the drift term, the dealer's central reservation price takes the explicit form 
\begin{align*}
    r(s, q, t) = s + \mu(T-t) - q\gamma\sigma^2(T-t).
\end{align*}
This formulation highlights two fundamental temporal dynamics: the directional anticipation and the horizon effect. While the $\mu(T-t)$ term adjusts the valuation based on the expected short-term trend, the risk premium associated with the inventory ($q\gamma\sigma^2(T-t)$) decays linearly as the terminal time $T$ approaches. Intuitively, as the time-to-horizon $(T-t)$ diminishes, there is a lower probability that the mid-price will undergo significant adverse fluctuations before the final liquidation. Consequently, as $t$ approaches $T$, the indifference price $r(s, q, t)$ converges toward the mid-price $s$, and the dealer's quotes naturally realign with those of a symmetric strategy.

The primary objective of this experiment is to empirically measure the extent to which the inventory-aware strategy reduces the dispersion of terminal wealth compared to a naive symmetric benchmark. By maintaining a constant order-flow specification and fixed accounting metrics, the simulation tests whether active quote skewing stabilizes profits by reducing P\&L variance. Furthermore, we expect to observe an improvement in risk-adjusted returns, as evidenced by higher Sharpe-like ratios, and a more disciplined control over inventory excursions. Without an additional terminal liquidation penalty, the strategy should be interpreted as reducing inventory excursions relative to the naive strategy, not as guaranteeing convergence of $q_t$ to zero before the end of the horizon.

\subsection{Case 2: Geometric Brownian Motion with Finite Horizon}\label{sec:gbm}

The second benchmark replaces the Arithmetic Brownian Motion with a Geometric Brownian Motion (GBM). In this setting, the dynamics of the mid-price are governed by the following stochastic differential equation
\begin{align*}
    dS_t = \mu S_t dt + \sigma S_t dW_t.
\end{align*}
This specification addresses a primary limitation of the ABM by preserving the strict positivity of asset prices. More importantly, the transition to GBM introduces a multiplicative price risk structure where the instantaneous volatility is proportional to the current price level $S_t$. Consequently, the inventory risk becomes price-level dependent: the absolute variance of a position scales quadratically with the asset price, meaning that holding a fixed number of shares $q$ at a higher price exposes the dealer to significantly greater marked-to-market risk.

Under the Avellaneda--Stoikov framework, the dealer's objective with an inventory $q$ is to maximize the expected Constant Absolute Risk Aversion (CARA) utility of terminal wealth, $W_T = x + q S_T$, yielding the value function
\begin{align*}
    v(x, s, q, t) = \mathbb{E}_t \left[ -e^{-\gamma (x + q S_T)} \right].
\end{align*}
Unlike the ABM case, where $S_T$ is normally distributed and the expectation can be explicitly solved using the moment-generating function of a Gaussian variable, the GBM assumes $S_T$ is lognormally distributed. Specifically, evaluating the expected utility requires computing the expectation of an exponential function evaluated at a lognormal random variable
\begin{align*}
    \mathbb{E}_t \left[ -e^{-\gamma q S_T} \right] = \mathbb{E}_t \left[ -e^{-\gamma q s \exp\left((\mu - \frac{1}{2}\sigma^2)(T-t) + \sigma(W_T - W_t)\right)} \right].
\end{align*}
This lognormal expectation does not possess a simple closed-form expression analogous to the Gaussian ABM case, so the exact value function cannot be factored into a tractable exponential-affine form. Due to this structural hurdle, the GBM--CARA dealer problem does not yield the same simple exact closed-form reservation-price formula used in the ABM benchmark.

To obtain an implementable benchmark, we use a local Gaussian, or local variance, approximation. For a sufficiently short high-frequency horizon, we approximate the conditional moments of the lognormal price distribution by their first-order terms: the expected price is $\mathbb{E}_t[S_T] \approx s + \mu s(T-t)$, and the conditional variance is $\text{Var}_t(S_T) \approx (\sigma s)^2(T-t)$. By treating the price dynamics as locally Gaussian at each time step $t$, we define an effective absolute volatility $\sigma_{\text{eff}} = \sigma s$, which represents the instantaneous monetary fluctuations of the asset.

Substituting this state-dependent variance into the Gaussian moment-generating identity, the approximated value function becomes
\begin{align*}
    v(x, s, q, t) \approx -e^{-\gamma x} \exp \left( -\gamma q \big(s + \mu s (T-t)\big) + \frac{1}{2} \gamma^2 q^2 (\sigma s)^2 (T-t) \right).
\end{align*}
The reservation price $r(s, q, t)$ is defined via the indifference principle for acquiring one additional unit of inventory, such that $v(x - r(s,q,t), s, q+1, t) = v(x, s, q, t)$. Solving the approximate local-Gaussian algebraic equation yields the following heuristic reservation price. It should not be read as the exact Avellaneda--Stoikov solution under GBM--CARA utility
\begin{align*}
    r(s, q, t) \approx s + \mu s(T-t) - q\gamma(\sigma s)^2(T-t).
\end{align*}
This formula, utilized within our simulation engine, allows the agent to dynamically scale its risk premium in local monetary-risk units. As the price $s$ increases, the approximate term $q\gamma(\sigma s)^2(T-t)$ grows quadratically, so the market maker skews quotes more strongly for the same physical inventory. 

The primary objective of this simulation is to observe how this adaptive quoting behavior performs compared to a naive strategy that fails to account for the multiplicative nature of the risk. We expect the inventory-aware agent to maintain tighter control over its position limits during price rallies, effectively mitigating the "exploding" variance inherent in the GBM process. This case serves as a crucial bridge between the idealized linear world of the ABM and the more realistic, heteroskedastic dynamics of actual financial markets.

\subsection{Synthesis: ABM vs GBM}

Having defined both the Arithmetic (ABM) and Geometric (GBM) price dynamics within a finite-horizon framework, it is crucial to synthesize the theoretical divergence between the two models before examining the empirical results of our time-discretized Monte Carlo simulation. Although both models share the underlying mechanics of order arrival and dealer accounting, the nature of the stochastic price path fundamentally alters the risk landscape for the market maker.

The divergence is rooted in the structure of the variance. In the ABM framework (Section~\ref{sec:abm}), the inventory risk premium is strictly linear and price-independent: holding $q$ shares carries the same marked-to-market risk regardless of whether the asset trades at $\$10$ or $\$100$. This idealized, homoskedastic environment provides a clean baseline to observe the pure mechanics of the Avellaneda--Stoikov inventory control and the linear decay of the horizon effect as time $t$ approaches terminal time $T$.

Conversely, the transition to the GBM framework (Section~\ref{sec:gbm}) introduces heteroskedasticity into the price process. The instantaneous variance of the position scales quadratically with the asset's price level ($S_t^2$). Consequently, a dealer operating under GBM cannot rely on a static perception of risk. As the asset appreciates in value, the monetary impact of price fluctuations magnifies, demanding a highly adaptive quoting response. To mitigate this exploding variance, the inventory-aware agent utilizes a local variance approximation, substituting a constant $\sigma^2$ with an effective, state-dependent variance $(\sigma S_t)^2$.

However, despite this divergence in how inventory risk is priced into the reservation price, a structural similarity persists between the two frameworks: under the same exponential fill-intensity specification, the liquidity offset around the reservation price has the same functional form. The Avellaneda--Stoikov architecture separates inventory-risk management from the optimization of execution probability. Consequently, in this benchmark implementation, the distance between the dealer's quotes and the relevant reservation price is governed by the market's liquidity profile ($k$) and the dealer's risk aversion parameter ($\gamma$). The final quote levels are not identical across ABM and GBM, because the reservation price itself changes with the local monetary variance in the GBM approximation.

The following portfolio-risk section removes the single-asset restriction. Once the dealer quotes several correlated instruments, the same inventory-control logic must be expressed in vector form, because the risk of each quote depends not only on its own position but also on the covariance structure of the whole book.

\section{Multivariate Extension: Covariance-Aware Inventory Control}


\subsection{From Independent Copies to Correlated Multi-Asset Dynamics}

The single-asset Avellaneda--Stoikov model can be replicated independently across many instruments, but this does not produce a genuine multi-asset market maker. If prices are correlated, the risk of holding a vector of inventories is not the sum of isolated scalar risks. A long position in one asset can either amplify or hedge the risk created by another position. The practical consequence is that the dealer should not decide the bid and ask of asset $i$ using only $q_t^i$, but the quote should also react to the rest of the book.

This is the key modeling step behind covariance-aware market making. \textcite{CarteaJaimungalPenalva2015} develop the stochastic-control language used for market making, inventory penalties, and multi-instrument trading problems. \textcite{Gueant2016} presents the Avellaneda--Stoikov model together with multi-asset extensions, while \textcite{BergaultEvangelistaGueantVieira2022} studies closed-form approximations for multi-asset market making when direct dynamic programming becomes computationally expensive.

The present article is self-contained for the covariance-aware inventory formulas and the diagnostics reported below. The construction of multivariate normal shocks, Cholesky factorization, Brownian paths, and multivariate geometric Brownian motion is recalled here because those objects are needed to understand how covariance enters the inventory-control mechanism; the companion article by \textcite{RostagnoFlorio2026MonteCarlo} provides a broader Monte Carlo treatment.

Let the dealer quote $d$ assets. The mid-price and inventory are vectors:
\[
S_t=(S_t^1,\ldots,S_t^d)^\top,\qquad
q_t=(q_t^1,\ldots,q_t^d)^\top .
\]
To simulate correlated shocks, start from a covariance matrix
\[
\Sigma\in\mathbb{R}^{d\times d},
\]
assumed symmetric positive definite. Its Cholesky factorization is
\[
\Sigma=LL^\top,
\]
where $L$ is lower triangular. If
\[
Z_k\sim\mathcal{N}(0,I_d),
\]
then
\[
\varepsilon_k=\sqrt{\Delta t}\,LZ_k
\]
satisfies
\[
\mathbb{E}[\varepsilon_k]=0,
\qquad
\mathrm{Cov}(\varepsilon_k)=\Delta t\,\Sigma .
\]
Therefore, the transformation $LZ_k$ converts independent standard Gaussian shocks into correlated shocks with the desired instantaneous covariance structure.

For the arithmetic Brownian benchmark used in the covariance-aware inventory formulas, the vector price process is
\[
dS_t=\mu\,dt+L\,dW_t,
\qquad
\Sigma=LL^\top.
\]
On a discrete grid, this becomes
\[
S_{t_{k+1}}=S_{t_k}+\mu\,\Delta t+\sqrt{\Delta t}\,LZ_k,
\qquad
Z_k\sim\mathcal{N}(0,I_d).
\]
In the data-free benchmark used here, $\mu$ is set to zero. The objective is not to model directional alpha, but to isolate the effect of inventory control under correlated price risk.

The same Cholesky logic also applies to multivariate geometric Brownian motion. If $\Sigma$ denotes the instantaneous covariance matrix of log-returns, then
\[
\frac{dS_t^i}{S_t^i}
=
\mu_i\,dt+\sum_{j=1}^{d}L_{ij}\,dW_t^j,
\qquad i=1,\ldots,d,
\]
and the exact lognormal discretization is
\[
S_{t_{k+1}}^i
=
S_{t_k}^i
\exp\!\left[
\left(\mu_i-\frac{1}{2}\Sigma_{ii}\right)\Delta t
+
\sqrt{\Delta t}\,(LZ_k)_i
\right].
\]
Equivalently, in vector form,
\[
S_{t_{k+1}}
=
S_{t_k}\odot
\exp\!\left[
\left(\mu-\frac{1}{2}\operatorname{diag}(\Sigma)\right)\Delta t
+
\sqrt{\Delta t}\,LZ_k
\right],
\]
where $\odot$ denotes componentwise multiplication.

For the market-making formulas below, we keep the arithmetic Brownian notation because it gives a clean Gaussian expression for the marked-to-market inventory risk. In a geometric Brownian setting, the same idea can be applied locally by replacing the constant monetary covariance matrix $\Sigma$ with the state-dependent covariance
\[
\Sigma_S(t)
=
D(S_t)\,\Sigma\,D(S_t),
\qquad
D(S_t)=\operatorname{diag}(S_t^1,\ldots,S_t^d),
\]
which reflects the fact that GBM volatility is proportional to the current price level.

\subsection{Portfolio Risk, Covariance Pressure, and Marginal Quotes}

Suppose the inventory vector $q_t$ is held fixed over the remaining horizon
\[
\tau=T-t.
\]
Under the arithmetic Brownian benchmark,
\[
S_T-S_t\mid\mathcal{F}_t
\sim
\mathcal{N}(0,\tau\Sigma),
\]
so the marked-to-market exposure of the dealer is
\[
q_t^\top(S_T-S_t).
\]
Its conditional variance is
\[
\mathrm{Var}_t\!\left(q_t^\top(S_T-S_t)\right)
=
\tau\,q_t^\top\Sigma q_t.
\]

Under CARA preferences, the corresponding certainty-equivalent inventory charge is proportional to
\[
C_t(q_t)
=
\frac{\gamma}{2}\tau\,q_t^\top\Sigma q_t.
\]
The scalar single-asset penalty
\[
\frac{\gamma}{2}\tau\sigma^2q_t^2
\]
is therefore replaced by a quadratic portfolio penalty. Since $\Sigma$ is symmetric,
\[
\frac{\partial}{\partial q_i}\left(q_t^\top\Sigma q_t\right)
=
2(\Sigma q_t)_i,
\]
so the gradient of the certainty-equivalent penalty is
\[
\nabla_q C_t(q_t)
=
\gamma\tau\,\Sigma q_t.
\]
The vector $\Sigma q_t$ is the central object of the multivariate extension. It measures the covariance pressure generated by the whole inventory book, not only by the position in one asset.

The covariance-aware central reservation vector is
\[
m_t
=
S_t-\gamma\tau\,\Sigma q_t.
\]
Its $i$-th component depends on the whole inventory vector through $(\Sigma q_t)_i$. If asset $i$ is positively correlated with assets already held long, then $(\Sigma q_t)_i$ is positive and the quote center for asset $i$ moves downward: the dealer becomes less willing to buy more and more willing to sell. If asset $i$ hedges the book, the same term can be small or even negative, so the dealer may quote it more aggressively even when $q_t^i$ itself is close to zero.

For implementation, it is useful to distinguish the central reservation vector from the exact marginal bid and ask reservation prices. Let $e_i$ denote the $i$-th unit vector. The additional inventory charge created by buying one more unit of asset $i$ is
\[
C_t(q_t+e_i)-C_t(q_t)
=
\gamma\tau(\Sigma q_t)_i
+
\frac{\gamma}{2}\tau\Sigma_{ii}.
\]
The inventory charge released by selling one unit is
\[
C_t(q_t)-C_t(q_t-e_i)
=
\gamma\tau(\Sigma q_t)_i
-
\frac{\gamma}{2}\tau\Sigma_{ii}.
\]
Therefore, the marginal reservation prices used by the multi-asset code are
\[
r_{t,i}^{bid}
=
S_t^i-\gamma\tau(\Sigma q_t)_i
-
\frac{\gamma}{2}\tau\Sigma_{ii},
\]
and
\[
r_{t,i}^{ask}
=
S_t^i-\gamma\tau(\Sigma q_t)_i
+
\frac{\gamma}{2}\tau\Sigma_{ii}.
\]

With exponential fill intensities
\[
\lambda_i(\delta)=A_i e^{-k_i\delta},
\]
the liquidity half-spread for asset $i$ is
\[
\delta_i^*
=
\frac{1}{\gamma}
\log\!\left(1+\frac{\gamma}{k_i}\right).
\]
The resulting quotes are
\[
p_{t,i}^{bid}
=
r_{t,i}^{bid}-\delta_i^*,
\qquad
p_{t,i}^{ask}
=
r_{t,i}^{ask}+\delta_i^*.
\]
This formula preserves the Avellaneda--Stoikov separation between inventory risk and execution risk. The covariance pressure $\Sigma q_t$ determines the quote skew, the diagonal term $\Sigma_{ii}$ contributes the marginal inventory-risk adjustment, and $k_i$ controls the liquidity component of the spread.

The expression \emph{covariance pressure} may sound abstract, but the idea is simple. In a multi-asset book, the dealer should not only ask how many units of asset $i$ are currently held. The dealer should also ask how much asset $i$ contributes to the risk of the whole inventory book. Consider a two-asset inventory vector and a covariance matrix:
\[
q=
\begin{pmatrix}
q_1\\
q_2
\end{pmatrix},
\qquad
\Sigma=
\begin{pmatrix}
\sigma_1^2 & \sigma_{12}\\
\sigma_{12} & \sigma_2^2
\end{pmatrix}.
\]
The vector $\Sigma q$ measures how the current inventory interacts with the covariance structure of the assets. Its $i$-th component, $(\Sigma q)_i$, tells us how asset $i$ contributes to the risk of the whole book, not only to the risk of its own position.

A useful way to read this object is the following:
\[
(\Sigma q)_i
=
\operatorname{Cov}\!\left(\Delta S_i,\ q^\top \Delta S\right)
\quad
\text{per unit of time}.
\]
In words, $(\Sigma q)_i$ is the covariance between the price change of asset $i$ and the marked-to-market variation of the whole inventory book. Therefore, it tells the dealer whether buying or selling asset $i$ increases or reduces total inventory risk.

The practical interpretation is immediate. If $(\Sigma q)_i>0$, buying more of asset $i$ increases the risk of the book. The dealer should therefore move the reservation price downward, becoming less willing to buy and more willing to sell. If $(\Sigma q)_i<0$, buying asset $i$ can reduce the risk of the book, because that asset hedges the current inventory. If $(\Sigma q)_i\simeq 0$, asset $i$ has little marginal impact on total portfolio risk.

Hence, the difference between $q_i$ and $(\Sigma q)_i$ is crucial. The term $q_i$ tells the dealer how many units of asset $i$ are currently held. The term $(\Sigma q)_i$ tells the dealer how asset $i$ affects the risk of the entire book once correlations are taken into account.

\subsection{Two-Asset Inventory Pressure and Correlated Paths}

Suppose that the dealer is long 10 units of asset 1 and short 5 units of asset 2:
\[
q=
\begin{pmatrix}
10\\
-5
\end{pmatrix},
\qquad
\Sigma=
\begin{pmatrix}
0.04 & 0.03\\
0.03 & 0.09
\end{pmatrix}.
\]
The covariance pressure is
\[
\Sigma q
=
\begin{pmatrix}
0.04 & 0.03\\
0.03 & 0.09
\end{pmatrix}
\begin{pmatrix}
10\\
-5
\end{pmatrix}
=
\begin{pmatrix}
0.25\\
-0.15
\end{pmatrix}.
\]
The first component is positive. This means that buying more of asset 1 would increase the risk of the current book. The dealer should therefore quote asset 1 more defensively: lower bid, more incentive to sell.

The second component is negative. This means that buying asset 2 would reduce the risk of the current book. This happens because the dealer is short asset 2, while asset 2 is positively correlated with asset 1. Buying asset 2 partially hedges the existing long exposure in asset 1.

The total instantaneous inventory risk is
\[
q^\top\Sigma q
=
\begin{pmatrix}
10 & -5
\end{pmatrix}
\begin{pmatrix}
0.25\\
-0.15
\end{pmatrix}
=
3.25.
\]
Therefore, over a remaining horizon $\tau=T-t$, the conditional variance of the marked-to-market inventory exposure is $\tau q^\top\Sigma q=3.25\tau$.

This example explains why the multi-asset model cannot use $q_i$ alone. The dealer may be short asset 2, but whether that short position is dangerous or useful depends on how asset 2 co-moves with the rest of the book.

The same covariance idea appears when simulating correlated price paths. The goal is to start from independent Gaussian shocks and transform them into correlated shocks with the desired covariance matrix. Let
\[
Z_k=
\begin{pmatrix}
Z_{1,k}\\
Z_{2,k}
\end{pmatrix}
\sim
\mathcal{N}(0,I_2),
\qquad
\Sigma=
\begin{pmatrix}
0.04 & 0.03\\
0.03 & 0.09
\end{pmatrix}.
\]
The Cholesky factorization gives a lower triangular matrix $L$ such that $\Sigma=LL^\top$. In this example,
\[
L=
\begin{pmatrix}
0.20 & 0\\
0.15 & 0.2598
\end{pmatrix},
\qquad
LL^\top=
\begin{pmatrix}
0.04 & 0.03\\
0.03 & 0.09
\end{pmatrix}.
\]
Thus, if $Z_k$ has independent standard normal components, then $LZ_k$ has covariance matrix $\Sigma$. This is the key role of the Cholesky factor: it converts independent shocks into correlated shocks.

\paragraph{Multivariate arithmetic Brownian motion.}
For a multivariate arithmetic Brownian motion, price shocks enter additively:
\[
dS_t=\mu\,dt+L\,dW_t,
\qquad
S_{t_{k+1}}=S_{t_k}+\mu\Delta t+\sqrt{\Delta t}\,LZ_k.
\]
Assume zero drift, initial prices $S_0=(100,50)^\top$, and one simulated Gaussian draw
\[
\mu=
\begin{pmatrix}
0\\
0
\end{pmatrix},
\qquad
S_0=
\begin{pmatrix}
100\\
50
\end{pmatrix},
\qquad
Z_k=
\begin{pmatrix}
0.50\\
-1.00
\end{pmatrix}.
\]
The correlated shock is
\[
LZ_k
=
\begin{pmatrix}
0.20 & 0\\
0.15 & 0.2598
\end{pmatrix}
\begin{pmatrix}
0.50\\
-1.00
\end{pmatrix}
=
\begin{pmatrix}
0.10\\
-0.1848
\end{pmatrix}.
\]
Therefore, for a generic step $\Delta t$,
\[
S_{t_{k+1}}
=
S_{t_k}
+
\sqrt{\Delta t}
\begin{pmatrix}
0.10\\
-0.1848
\end{pmatrix}.
\]
In the arithmetic Brownian case, the shock is added directly to the price vector. The covariance matrix $\Sigma$ is therefore interpreted as the covariance matrix of price changes.

\paragraph{Multivariate geometric Brownian motion.}
For a multivariate geometric Brownian motion, the same Cholesky idea is used, but the shocks enter multiplicatively. In this case, the covariance matrix refers to log-returns rather than price changes. Let
\[
\Sigma_{\log}
=
\begin{pmatrix}
0.04 & 0.03\\
0.03 & 0.09
\end{pmatrix}.
\]
This means that asset 1 has volatility $20\%$, asset 2 has volatility $30\%$, and their log-return covariance is $0.03$.

The dynamics are
\[
\frac{dS_t^i}{S_t^i}
=
\mu_i\,dt+\sum_{j=1}^{2}L_{ij}\,dW_t^j,
\qquad i=1,2.
\]
The exact discrete-time simulation is
\[
S_{t_{k+1}}^i
=
S_{t_k}^i
\exp\!\left[
\left(\mu_i-\frac{1}{2}\Sigma_{\log,ii}\right)\Delta t
+
\sqrt{\Delta t}(LZ_k)_i
\right].
\]
Equivalently, in vector form,
\[
S_{t_{k+1}}
=
S_{t_k}
\odot
\exp\!\left[
\left(\mu-\frac{1}{2}\operatorname{diag}(\Sigma_{\log})\right)\Delta t
+
\sqrt{\Delta t}LZ_k
\right],
\]
where $\odot$ denotes componentwise multiplication.

Using again the same Gaussian draw,
\[
Z_k=
\begin{pmatrix}
0.50\\
-1.00
\end{pmatrix},
\qquad
LZ_k=
\begin{pmatrix}
0.10\\
-0.1848
\end{pmatrix},
\]
and assuming zero drift, the two simulated prices become
\[
S_{t_{k+1}}^1
=
S_{t_k}^1
\exp\!\left[
-\frac{1}{2}(0.04)\Delta t
+
0.10\sqrt{\Delta t}
\right],
\qquad
S_{t_{k+1}}^2
=
S_{t_k}^2
\exp\!\left[
-\frac{1}{2}(0.09)\Delta t
-
0.1848\sqrt{\Delta t}
\right].
\]
The correlation mechanism is therefore the same as in the arithmetic Brownian case: Cholesky is used to create correlated Gaussian shocks. The difference is how those shocks affect prices. In arithmetic Brownian motion, shocks are added to prices. In geometric Brownian motion, shocks are added to log-prices and therefore affect prices multiplicatively.

\subsection{GBM Inventory Penalty and Implementation Design}

The distinction between arithmetic and geometric Brownian motion matters for inventory control. In the arithmetic Brownian model, $\Sigma$ is already the covariance matrix of price changes, so the inventory risk penalty can directly use $q^\top\Sigma q$.

In the geometric Brownian model, however, $\Sigma_{\log}$ is the covariance matrix of log-returns. If inventories are measured in units of shares or contracts, the dealer needs the covariance matrix of price changes. Locally, as a state-dependent approximation rather than as a closed-form HJB solution, this is approximated by
\[
\Sigma_S(t)
=
D(S_t)\Sigma_{\log}D(S_t),
\qquad
D(S_t)=\operatorname{diag}(S_t^1,S_t^2).
\]
For example, if $S_t=(100,50)^\top$, then
\[
D(S_t)=
\begin{pmatrix}
100 & 0\\
0 & 50
\end{pmatrix},
\qquad
\Sigma_S(t)
=
\begin{pmatrix}
400 & 150\\
150 & 225
\end{pmatrix}.
\]
Using the same inventory vector as before,
\[
q=
\begin{pmatrix}
10\\
-5
\end{pmatrix},
\qquad
\Sigma_S(t)q
=
\begin{pmatrix}
400 & 150\\
150 & 225
\end{pmatrix}
\begin{pmatrix}
10\\
-5
\end{pmatrix}
=
\begin{pmatrix}
3250\\
375
\end{pmatrix}.
\]
This final calculation highlights an important practical point. In a GBM setting, the risk of holding one unit of an expensive asset is mechanically larger than the risk of holding one unit of a cheap asset, because volatility scales with the price level. Therefore, when the inventory is measured in physical units, the inventory penalty should use the local price covariance matrix $\Sigma_S(t)$ rather than the log-return covariance matrix $\Sigma_{\log}$ directly.

The illustrative calculations above use two assets because the covariance-pressure mechanism can be seen directly in a $2\times 2$ matrix. The numerical results reported later use a three-asset diagnostic. In both cases, at each time step, correlated Gaussian shocks are generated through the Cholesky factor $L$ of $\Sigma$. The same shock matrix is reused for the naive and inventory-aware policies. The same matrix of uniform random variables is also reused to simulate fills. This common-random-number design ensures that differences in terminal wealth and inventory are driven by the quoting rule, not by different simulated scenarios.

The execution model is intentionally simple. Conditional on the quoted distances, fills are independent across assets, and each asset has its own parameters $A_i$ and $k_i$. This isolates the covariance channel in the inventory term. More realistic features such as cross-excitation of order flow, queue priority, market impact, client segmentation, latency, and variable trade sizes are important extensions, but they would introduce additional microstructure mechanisms that are outside the purpose of the present benchmark.

The multivariate output reports terminal PnL, terminal inventory vectors, the risk norm
\[
\sqrt{q_T^\top\Sigma q_T},
\]
and the covariance pressure
\[
\Sigma q_T.
\]
These diagnostics answer a different question from the scalar Monte Carlo table. The issue is not only whether each individual inventory is smaller, but whether the remaining book is less risky after accounting for the correlation structure of the assets.

Chapter~\ref{ch:numerics} turns these models into a reproducible experimental design. The same logic is implemented through a common simulation engine so that the single-asset and multi-asset cases can be evaluated under comparable accounting, fill, and Monte Carlo conventions.

The covariance-aware pressure $\Sigma q$ already shows why separate inventories can be an inadequate state for portfolio risk. The memory extension asks a different question: whether current volatility is enough to predict future $\Sigma$. This motivates causal, history-dependent drivers without changing the meaning of a fill, cash or inventory.
